\clearpage
\fancyhead[L]{\small GAUSS AN BODE.}
\hypertarget{native-283-lower}{}
\begin{center}
GAUSS AN BODE.
\end{center}

\begin{center}
Astronomisches Jahrbuch für 1809. S. 215–219. Berlin 1806.
\end{center}

\begin{flushright}
Braunschweig, 30. Juli 1806.
\end{flushright}

\par\indent In meinem letzten Brief habe ich Ew. --- versprochen, Ihnen die Ergebnisse meiner neuesten Rechnungen über die neuen Planeten mitzuteilen, sobald ich damit fertig sei. Ich beeile mich nun, dieses Versprechen einzulösen.

\par\indent Meine Beobachtungen der Pallas vom Februar dieses Jahres habe ich Ihnen bereits mitgeteilt (siehe \href{https://www.digitale-sammlungen.de/en/view/bsb10538338?page=147}{S.~137} [\hyperlink{native-272}{S.~267} dieses Bandes]). Für die letzte Verbesserung der Elemente habe ich jedoch nur die in Mailand gemachten Beobachtungen verwendet, die inzwischen auch im \href{https://www.digitale-sammlungen.de/en/view/bsb10538606?page=100}{Juliheft der M.\,C.} abgedruckt sind. Dies sind demnach meine neuesten elliptischen Elemente der Pallas.

\begin{center}
\emph{Elemente der Pallas.} (IX.)

Epochen der Länge, Meridian von Seeberg.
\end{center}

\begin{center}\begingroup
\fontsize{9}{11}\selectfont
\setlength{\tabcolsep}{3pt}\renewcommand{\arraystretch}{1.12}
\begin{tabular}{ll}
1802.  & $143^\circ\,24'\,1''1$ \\
1803.  & $221\;32\;56.4$ \\
1804.  & $299\;54\;42.5$ \\
1805.  & $18\;3\;37.8$ \\
1806.  & $96\;12\;33.2$ \\
1807.  & $174\;21\;28.5$ \\
\end{tabular}\endgroup\end{center}
\begin{center}\begingroup
\fontsize{9}{11}\selectfont
\setlength{\tabcolsep}{3pt}\renewcommand{\arraystretch}{1.12}
\begin{tabular}{ll}
tägliche mittlere Bewegung .\ .\ .\ .\ . & $770''7816$ \\
Aphel 1806 .\ .\ .\ .\ .\ .\ .\ . & $301^\circ\,12'\,21''$ \\
aufsteigender Knoten 1806 .\ .\ .\ . & $172\;30\;56$ \\
Neigung der Bahn .\ .\ .\ .\ .\ .\ . & $34\;37\;8$ \\
Exzentrizität .\ .\ .\ .\ .\ .\ .\ .\ .\ . & $0.245384$ \\
Logarithmus der Halbachse .\ .\ .\ .\ . & $0.4421021$ \\
\end{tabular}\endgroup\end{center}
\SourceLeaf{284}%
\par\indent Nach diesen Elementen habe ich den Lauf der Pallas für ihre nächste Erscheinung berechnet. Ich habe die Ehre, die Ephemeride hier beizufügen. Prof. \textsc{Harding} hat wieder eine Karte dazu angefertigt, die Herr \textsc{v. Zach} für die M.\,C. stechen lassen wird. Von Prof. \textsc{Harding}s Atlas, der für die Beobachtung der Asteroiden so notwendig ist, werden schon demnächst ein paar Blätter erscheinen.

\par\indent Auch die hier gemachten Beobachtungen der Juno habe ich Ihnen bereits mitgeteilt (siehe \href{https://www.digitale-sammlungen.de/en/view/bsb10538338?page=147}{S.~137} [\hyperlink{native-272}{S.~267} dieses Bandes]). Soweit ich bisher sehe, waren es die ersten in diesem Jahr. Auch hier habe ich zur Verbesserung der Elemente hauptsächlich die Beobachtungen von \textsc{Oriani} verwendet, die im \href{https://www.digitale-sammlungen.de/en/view/bsb10538606?page=98}{Juliheft der M.\,C.} abgedruckt sind.

\begin{center}
VI. \emph{Elemente der Juno.}

Epochen der mittleren Länge. Meridian von Seeberg.
\end{center}

\begin{center}\begingroup
\fontsize{9}{11}\selectfont
\setlength{\tabcolsep}{3pt}\renewcommand{\arraystretch}{1.12}
\begin{tabular}{ll}
1805.  & $42^\circ\,35'\,7''3.$ \\
1806.  & $125^\circ\,11'\,20''1.$ \\
1807.  & $207^\circ\,47'\,33''0$ \\
\end{tabular}\endgroup\end{center}

\begin{center}\begingroup
\fontsize{9}{11}\selectfont
\setlength{\tabcolsep}{3pt}\renewcommand{\arraystretch}{1.12}
\begin{tabular}{ll}
Länge des Aphels 1806 .\ .\ . & $233^\circ\,17'\,1''$ \\
Länge des aufsteigenden Knotens 1806  & $171\;4\;57$ \\
Neigung der Bahn .\ .\ .\ .\ .\ .\ . & $13\;3\;28''4$ \\
tägliche mittlere tropische Bewegung . & $814''720$ \\
Exzentrizität .\ .\ .\ .\ .\ .\ .\ .\ . & $0.2549441$ \\
Logarithmus der Halbachse .\ .\ .\ . & $0.4260480$ \\
\end{tabular}\endgroup\end{center}

\par\indent Die nach diesen Elementen berechnete Ephemeride habe ich auch schon Herrn Prof. \textsc{Harding} mitgeteilt, damit er danach die Karte zeichnen kann. Sie erhalten hier eine Abschrift davon. Im nächsten Jahr wird es jedoch äußerst schwierig sein, Juno zu beobachten. Schon in diesem Jahr war sie äußerst lichtschwach, und im nächsten wird sie nicht einmal halb so hell erscheinen wie in diesem\HistoricalNote{24}. [Der Lauf der Juno vom 7. Dezember 1806 bis zum 21. September 1807 ist in Abständen von drei Tagen für 12 Uhr in Seeberg angegeben, mit Rektaszension und Deklination auf die Minute genau.] Die $\text{\AstroSymbols ☍}$ der $\PallasSymbol$ tritt demnach am 4. Mai 1807 um 21 Uhr bei $223^\circ\,53'$ Länge und $42^\circ\,18'$ nördlicher Breite ein.

\par\indent Über Ceres habe ich dieses Jahr noch keine weiteren Rechnungen angestellt. Mir fehlen noch gute Beobachtungen in der Nähe der Opposition. Da die $\CeresSymbol$ eine so günstige Lage hatte, zweifle ich nicht daran, dass diese $\text{\AstroSymbols ☍}$ an mehreren Orten gut beobachtet worden sein wird. Sollten Sie solche Beobachtungen haben, wäre es mir eine Gefälligkeit, wenn Sie sie mir \emph{bald} mitteilen. Gern werde ich Ihnen dann auch die Ergebnisse daraus zuschicken\renewcommand{\thefootnote}{**)}\HistoricalNote{25}\renewcommand{\thefootnote}{*)}.



